\documentclass[11pt,a4paper]{article}

\usepackage[T1]{fontenc}
\usepackage[utf8]{inputenc}
\usepackage{lmodern}
\usepackage{amssymb,amsmath,amsthm,mathtools}
\usepackage[margin=2.5cm]{geometry}
\usepackage{parskip}
\usepackage[hidelinks]{hyperref}
\usepackage{microtype}

% Theorem environments (arXiv / journal style)
\theoremstyle{plain}
\newtheorem{theorem}{Theorem}[section]
\newtheorem{proposition}[theorem]{Proposition}
\newtheorem{lemma}[theorem]{Lemma}
\newtheorem{corollary}[theorem]{Corollary}

\theoremstyle{definition}
\newtheorem{definition}[theorem]{Definition}
\newtheorem{example}[theorem]{Example}
\newtheorem{remark}[theorem]{Remark}

% Notation
\newcommand{\Ext}{\mathrm{Ext}}
\newcommand{\Shape}{\mathrm{Shape}}
\newcommand{\Pos}{\mathrm{Pos}}
\newcommand{\Type}{\mathbf{Type}}
\newcommand{\id}{\mathrm{id}}
\newcommand{\Cont}{\mathbf{Cont}}
\newcommand{\seqop}{\mathbin{\triangleleft}}

\title{The Sequential Operator on Containers:\\
  Units, Associator, and Composition of Extensions}
\author{Extracted from a machine-checked development}
\date{}

\begin{document}
\maketitle

\begin{abstract}
We equip the category of containers with the \emph{sequential operator}
$\seqop$, whose extension is composition of the induced polynomial
functors: $\Ext(G \seqop F) \cong \Ext\,G \circ \Ext\,F$. The unit is the
identity container $I = 1 \seqop {!}$. We construct the left and right
unitors $I \seqop F \cong F$ and $F \seqop I \cong F$, and the
associator $(H \seqop G) \seqop F \cong H \seqop (G \seqop F)$, as
isomorphisms of containers. All dependent types involved agree
definitionally: the associator is re-association of an iterated
dependent sum, and no transport is required. A $\seqop$-comonoid in the
category of containers is precisely a directed container. All stated
results have been machine-checked.
\end{abstract}

\tableofcontents
\bigskip

%======================================================================
\section{Introduction}
%======================================================================

Containers present a class of polynomial endofunctors on types: a shape
type equipped with a family of position types determines the functor
\[
  X \;\longmapsto\; \sum_{s}\; X^{P(s)}.
\]
Composition of such functors is again a polynomial functor of the same
form, and the composite has an explicit combinatorial description. That
description is the sequential operator~$\seqop$.

The operator is the spine of the theory of directed containers: a
directed container is exactly a $\seqop$-comonoid in the category of
containers. The present note records the monoidal structure itself---the
operator, its unit, and the coherence isomorphisms---leaving the
comonoid laws to the directed-container development.

\paragraph{What this note claims.}
In ordinary mathematical language: the sequential composite
$G \seqop F$ has shapes
\[
  \sum_{t \,:\, \Shape(G)}\;
  \bigl(\Pos(G)(t) \to \Shape(F)\bigr)
\]
and positions that are pairs of an outer $G$-position with an inner
$F$-position in the shape sitting there. Its extension is functor
composition. The monoidal coherences (unitors and associator) are
isomorphisms of containers whose underlying shape maps are the evident
bijections, and all position fibres line up definitionally.

%======================================================================
\section{Containers recalled}
%======================================================================

\begin{definition}[Container]\label{def:container}
A \textbf{container} $C$ consists of a type $\Shape(C)$ of shapes and,
for each shape~$s$, a type $\Pos(C)(s)$ of positions. We write
$C = (S, P)$ when convenient.
\end{definition}

\begin{definition}[Extension]\label{def:ext}
The \textbf{extension} of a container~$C$ at a type~$X$ is
\begin{equation}
  \Ext\,C\,X \;=\; \sum_{s \,:\, \Shape(C)}\;
  \bigl(\Pos(C)(s) \to X\bigr).
\end{equation}
\end{definition}

\begin{definition}[Container morphism]\label{def:morphism}
A \textbf{morphism} $\varphi : C \Rightarrow D$ consists of a map on
shapes $u : \Shape(C) \to \Shape(D)$ and, for each shape~$s$, a
contravariant map on positions
$g_s : \Pos(D)(u(s)) \to \Pos(C)(s)$.
\end{definition}

We write $\Cont$ for the category of containers and container morphisms.
Composition of morphisms and identities are the evident ones
(covariant on shapes, contravariant on positions).

%======================================================================
\section{The sequential operator}
%======================================================================

\begin{definition}[Sequential operator]\label{def:seq}
Let $F$ and $G$ be containers. The \textbf{sequential composite}
$G \seqop F$ (``$G$ outside, $F$ inside'') is the container with
\begin{align}
  \Shape(G \seqop F)
  &\;=\;
  \sum_{t \,:\, \Shape(G)}\;
  \bigl(\Pos(G)(t) \to \Shape(F)\bigr), \\[4pt]
  \Pos(G \seqop F)\,(t, f)
  &\;=\;
  \sum_{q \,:\, \Pos(G)(t)}\;
  \Pos(F)\bigl(f(q)\bigr).
\end{align}
A shape is an outer $G$-shape together with a choice of inner $F$-shape
for each $G$-position; a position is an outer position together with an
inner position in the shape sitting there.
\end{definition}

\begin{definition}[Identity container]\label{def:I}
The \textbf{identity container} is
\begin{equation}
  I \;=\; \bigl(1,\; {\ast} \mapsto 1\bigr),
\end{equation}
one shape and one position. Its extension is the identity functor, and
it is the unit for~$\seqop$.
\end{definition}

%======================================================================
\section{Extension is functor composition}
%======================================================================

The defining property of~$\seqop$ is that extension turns sequential
composition into composition of functors.

\begin{definition}[Uncurrying of labellings]\label{def:seqext}
For containers $G$ and~$F$ and a type~$X$, define maps
\begin{align}
  \alpha_X &\;:\;
  \Ext\,(G \seqop F)\,X
  \;\longrightarrow\;
  \Ext\,G\,(\Ext\,F\,X), \\[4pt]
  \beta_X &\;:\;
  \Ext\,G\,(\Ext\,F\,X)
  \;\longrightarrow\;
  \Ext\,(G \seqop F)\,X
\end{align}
by regrouping: $\alpha$ sends a labelling of the nested position sum to
an outer structure of inner structures; $\beta$ is the inverse
uncurrying. Explicitly, writing a shape of $G \seqop F$ as
$\langle t, f \rangle$ and a labelling as~$v$,
\begin{align}
  \alpha_X\,\bigl\langle \langle t, f \rangle,\, v \bigr\rangle
  &\;=\;
  \bigl\langle
    t,\;
    q \mapsto
    \bigl\langle f(q),\; r \mapsto v\langle q, r \rangle \bigr\rangle
  \bigr\rangle, \\[4pt]
  \beta_X\,\bigl\langle t,\, m \bigr\rangle
  &\;=\;
  \bigl\langle
    \langle t,\, q \mapsto \pi_1(m(q)) \rangle,\;
    \langle q, r \rangle \mapsto \pi_2(m(q))(r)
  \bigr\rangle.
\end{align}
\end{definition}

\begin{theorem}[Extension of sequential composite]
\label{thm:seq-ext}
For all containers $G$ and~$F$ and every type~$X$, the maps $\alpha_X$
and $\beta_X$ are mutually inverse:
\begin{equation}
  \beta_X \circ \alpha_X \;=\; \id,
  \qquad
  \alpha_X \circ \beta_X \;=\; \id.
\end{equation}
Thus $\Ext\,(G \seqop F) \cong \Ext\,G \circ \Ext\,F$ as functors on
types. Both fibres agree definitionally, so the round trips hold after
case analysis on nested pairs.
\end{theorem}

\begin{proof}
By lean.
\end{proof}

%======================================================================
\section{Isomorphisms of containers}
%======================================================================

\begin{definition}[Container isomorphism]\label{def:iso}
An \textbf{isomorphism of containers} $C \cong D$ is a pair of
morphisms $\mathrm{hom} : C \Rightarrow D$ and
$\mathrm{inv} : D \Rightarrow C$ that are mutual inverses under
composition of container morphisms.
\end{definition}

\begin{remark}
The unitors and associator are isomorphisms, not equalities: the
underlying shape types are bijective (e.g.\ $(1 \to S) \cong S$) but not
definitionally identical.
\end{remark}

%======================================================================
\section{Left and right unitors}
%======================================================================

\begin{theorem}[Left unitor]\label{thm:left-unitor}
For every container~$F$ there is an isomorphism of containers
\begin{equation}
  I \seqop F \;\cong\; F.
\end{equation}
Forward, read off the single inner shape
($f : 1 \to \Shape(F)$ maps to $f(\ast)$); backward, send a shape~$s$ to
the constant assignment ${\ast} \mapsto s$. Positions transport along the
same collapse of the unit type. Both round trips hold by $\eta$-laws for
the unit type and dependent sums.
\end{theorem}

\begin{proof}
By lean.
\end{proof}

\begin{theorem}[Right unitor]\label{thm:right-unitor}
For every container~$F$ there is an isomorphism of containers
\begin{equation}
  F \seqop I \;\cong\; F.
\end{equation}
Forward, drop the trivial inner data
($g : \Pos(F)(s) \to 1$ is unique); backward, pad with the unique map
to the unit type. Positions
$\sum_{p \,:\, \Pos(F)(s)} 1$ collapse to $\Pos(F)(s)$.
\end{theorem}

\begin{proof}
By lean.
\end{proof}

%======================================================================
\section{The associator}
%======================================================================

\begin{theorem}[Associator]\label{thm:associator}
For all containers $H$, $G$, and~$F$ there is an isomorphism of
containers
\begin{equation}
  (H \seqop G) \seqop F
  \;\cong\;
  H \seqop (G \seqop F).
\end{equation}
\end{theorem}

\begin{proof}
By lean.
\end{proof}

\begin{remark}[Explicit data]
A shape of the left-hand side is a triple
$\bigl(\langle a, g \rangle,\, k\bigr)$ with
$g : \Pos(H)(a) \to \Shape(G)$ and
\[
  k \;:\;
  \Bigl(
    \sum_{u \,:\, \Pos(H)(a)}\;
    \Pos(G)\bigl(g(u)\bigr)
  \Bigr)
  \to
  \Shape(F).
\]
A shape of the right-hand side is a pair $\langle a, m \rangle$ with
\[
  m \;:\;
  \Pos(H)(a)
  \to
  \sum_{b \,:\, \Shape(G)}\;
  \bigl(\Pos(G)(b) \to \Shape(F)\bigr).
\]
The bijection curries the inner data:
\begin{equation}
  m(u) \;=\; \bigl\langle g(u),\; v \mapsto k\langle u, v \rangle \bigr\rangle,
\end{equation}
inverted by $g(u) = \pi_1(m(u))$ and
$k\langle u, v \rangle = \pi_2(m(u))(v)$.

On positions the isomorphism is re-association of a nested dependent
sum,
\begin{equation}
  \bigl\langle \langle u, v \rangle,\, p \bigr\rangle
  \;\longleftrightarrow\;
  \bigl\langle u,\, \langle v, p \rangle \bigr\rangle.
\end{equation}
 Crucially, the fibre
$\Pos(F)\bigl(\pi_2(m(u))(v)\bigr)$ is definitionally equal to
$\Pos(F)\bigl(k\langle u, v \rangle\bigr)$, so no transport is needed
and the round trips close by case analysis on nested pairs.
\end{remark}

%======================================================================
\section{Organising remark: comonoids and directed containers}
%======================================================================

\begin{remark}[$\seqop$-comonoids]
A \textbf{$\seqop$-comonoid} in $\Cont$ is a container~$C$ equipped with
a counit $C \Rightarrow I$ and a comultiplication
$C \Rightarrow C \seqop C$ satisfying the comonoid laws (coassociativity
up to the associator; counitality up to the unitors). Such a structure is
precisely a \textbf{directed container}: equivalently, a comonad on
types whose underlying functor is a container extension.

Thus $(\Cont, \seqop, I)$ is the monoidal home of directed containers.
The comonoid laws unfold to the directed-container axioms (root,
subtree, shift, and laws D1--D5); we record the correspondence as the
organising remark and leave the equivalence to the directed-container
development.
\end{remark}

%======================================================================
\section{Discussion}
%======================================================================

\paragraph{What has been established.}
The sequential operator equips containers with a monoidal structure
whose extension is functor composition. The unitors and associator are
isomorphisms of containers; all relevant fibres agree definitionally, so
the constructions are transport-free.

\paragraph{Why definitional agreement matters.}
In dependent type theory, re-associating nested sums often forces
transport along equalities of fibres. Here the position type of
$G \seqop F$ is arranged so that the two associations of a triple sum
are definitionally the same family. The associator is therefore nothing
but the combinatorial re-bracketing
$((u,v),p) \leftrightarrow (u,(v,p))$.

\paragraph{Machine checking.}
Every definition and theorem statement above corresponds to a declaration
in a machine-checked formalisation. The phrase ``by lean'' indicates that
the proof has been verified by a proof assistant.

%======================================================================
\section*{Acknowledgements}
%======================================================================

The mathematical content follows the theory of containers of Abbott,
Altenkirch, and Ghani, and the sequential (composition) product standard
in that literature. The presentation here is a prose rendering of a
formal development intended for human review of the theorem statements.

\begin{thebibliography}{9}

\bibitem{abbott2005}
M.~Abbott, T.~Altenkirch, and N.~Ghani.
\newblock Containers: constructing strictly positive types.
\newblock \emph{Theoretical Computer Science}, 342(1):3--27, 2005.

\bibitem{abbott2003}
M.~Abbott, T.~Altenkirch, and N.~Ghani.
\newblock Categories of containers.
\newblock In \emph{FoSSaCS}, volume~2620 of \emph{LNCS}, pages 23--38.
\newblock Springer, 2003.

\bibitem{ahman2012}
D.~Ahman, J.~Chapman, and T.~Uustalu.
\newblock When is a container a comonad?
\newblock In \emph{FoSSaCS}, volume~7213 of \emph{LNCS}, pages 74--88.
\newblock Springer, 2012.

\end{thebibliography}

\end{document}
